\section{CS224N 的主线：我们到底学到了什么}

\subsection{从 word vectors 到 foundation models}

如果把整门课压缩成一句话，那就是：我们从“如何给词建一个几何空间”出发，最后走到了“如何让模型成为一个可以做事的通用语言系统”。

\begin{enumerate}
    \item \textbf{Word vectors}：词的意义可以由其上下文来定义。
    \item \textbf{Neural NLP}：用可学习的表示替代手工特征。
    \item \textbf{Sequence models}：RNN、LSTM 处理序列依赖。
    \item \textbf{Transformers}：通过 attention 和并行化把大模型训练变成现实。
    \item \textbf{Pretraining + post-training}：大规模预训练再做任务对齐，形成 foundation model。
    \item \textbf{Evaluation and reasoning}：模型越来越强，但评测、推理和可靠性仍然是难点。
\end{enumerate}

\begin{figure}[H]
\centering
\includegraphics[width=0.95\textwidth]{slides-images/slide-020.jpg}
\caption{GPT-4 在真实工作场景中带来生产力提升，但效果显著依赖任务与使用方式\protect\footnotemark}
\end{figure}
\footnotetext{Slides 中展示了咨询任务里的结果分布。}

\begin{warningbox}{“看起来很强”不等于“已经理解”}
LLM 在很多任务上确实很有用，但这不自动说明它们已经具备和人类相同的泛化能力。更现实的说法是：模型已经变得非常会“利用经验”，但是否真正“理解”语言，仍然是开放问题。
\end{warningbox}

\subsection{本章小结}

课程前半段解决的是“怎么建模”，后半段逐步转向“这些模型到底在学什么、能做什么、还缺什么”。最后一讲把这两类问题放到一起：技术进步是真实的，但它并没有关闭更深层的科学问题。

%% ========== 开放问题 ==========

